\documentclass{oxmathproblems}

\printanswers

\course{Understanding Analysis, 2nd Edition: Stephen Abbott}
\sheettitle{Chapter 4 - Section 4.2 - Functional Limits} 

\begin{document}
\begin{questions}

%------------------------- Problem 1 -------------------------
\miquestion
\begin{enumerate}[label=(\alph*)]
  \item Supply the details of how Corollary 4.2.4 part (ii) follows from the Sequential Criterion for Functional Limits in Theorem 4.2.3 and the Algebraic Limit Theorem
  for sequences proved in Chapter 2.
  \item Now, write another proof of Corollary 4.2.4 part (ii) directly from Definition 4.2.1 without using the sequential criterion.
  \item Repeat (a) and (b) for Corollary 4.2.4 part (iii)
\end{enumerate}

\begin{solution}
  \begin{enumerate}[label=(\alph*)]
    \item We need to prove that $\lim_{x \to c} [f(x) + g(x)] = L + M$ given that the Sequential Criterion holds for $f$ and $g$.
    
    By the Criterion, given any $(x_{n}) \to c$, we have $f(x_{n}) \to L$ and $g(x_{n}) \to M$, so by the Algebraic Limit Theorem
    $f(x_{n}) + g(x_{n}) \to L + M$. This shows that $f+g$ satisfies the Sequential Criterion. (By the way, just a reminder about the possibly ambiguous notation: $f(x_{n})$ is the sequence
    $f(x_{1}), f(x_{2}), f(x_{3}), ...$ and \textit{not} $f(x_{1}, x_{2}, x_{3}, ...)$.

    By Theorem 4.2.3, if $f + g$ satisfies the Sequential Criterion, it also satisfies $\lim_{x \to c} [f(x) + g(x)] = L + M$.

    \item There exists $\delta_{f} > 0$ such that $|x-c| < \delta_{f} \implies |f(x)-L| < \frac{\epsilon}{2}$. Similarly choose $\delta_{g}$
    such that $|x-c| < \delta_{g} \implies |g(x) - M| < \frac{\epsilon}{2}$. Now let $\delta = \min(\delta_{f}, \delta_{g})$, then
    $|x-c| < \delta \implies |f(x)+g(x) - (L+M)| \leq |f(x)-L| + |g(x)-M| < \frac{\epsilon}{2} + \frac{\epsilon}{2} = \epsilon$.
  
    \item Part (a).
    
    Same thing. By the Algebraic Limit Theorem $f(x_{n})g(x_{n}) \to LM$ so $(fg)(x_{n})$ satisfies the Sequential Criterion. Then by Theorem 4.2.3
    $\lim_{x \to c}[f(x)g(x)] = LM$.

    Part (b).

    We need to prove that for any $\epsilon > 0$, there exists $\delta$ such that $|x-c| < \delta \implies |f(x)g(x) - LM| < \epsilon$.
    
    We know we'll need $|f(x) - L|$ and $|g(x) - M|$ to appear. Let's start with $|f(x) - L|$, what do we need to add to
    $f(x)g(x)$ to be able to factor out $f(x)-L$? We need to add $g(x)L$, which would get us:
    
    \begin{align*}
      |f(x)g(x) - LM| 
      & = |f(x)g(x) - g(x)L + g(x)L - LM| \\
      & = |g(x)(f(x)-L) + L(g(x)-M)| \\
      & \leq |g(x)||f(x)-L| + |L||g(x)-M|
    \end{align*}

    It turns out that we were lucky; the $g(x)L$ factor also happens to factor out nicely with $- LM$ to create the $|g(x)-M|$ we need.

    Now the problem is how to bound $|g(x)|$. We can pick $\delta_{g}$ to bound $|g(x)-M| < \frac{\epsilon}{2|L|}$,
    which will bound $|g(x)| < |M| + \frac{\epsilon}{2|L|}$. Pick $\delta_{f}$ such that $|f(x)-L| < \frac{\epsilon}{2(|M| + \frac{\epsilon}{2|L|})}$.
    Finally, pick $\delta = \min(\delta_{f}, \delta_{g})$, then:

    \begin{align*}
      |f(x)g(x) - LM|
      & \leq |g(x)||f(x)-L| + |L||g(x)-M| \\
      & < \left(|M| + \frac{\epsilon}{2|L|}\right) \frac{\epsilon}{2(|M| + \frac{\epsilon}{2|L|})} + |L|\frac{\epsilon}{2|L|} \\
      & = \epsilon
    \end{align*}
  \end{enumerate}

  But note that this fails for the case $L = 0$, because there'll be division by zero. We could treat this case separately, or we can choose a smarter bound that forces the
  denominators to be nonzero.

  Bound $|g(x)-M| < \frac{\epsilon}{2|L|+1}$. This clearly works, since the denominator is even bigger than before, and now it can't be zero. The bound for
  $g(x)$ and $|f(x)-L|$ changes accordingly, and we'll have:

  \begin{align*}
    |f(x)g(x) - LM|
    & \leq |g(x)||f(x)-L| + |L||g(x)-M| \\
    & < \left(|M| + \frac{\epsilon}{2|L|+1}\right) \frac{\epsilon}{2\left(|M| + \frac{\epsilon}{2|L|+1}\right)} + |L|\frac{\epsilon}{2|L|+1} \\
    & < \frac{\epsilon}{2} + \frac{\epsilon}{2} \\
    & = \epsilon
  \end{align*}
\end{solution}


%------------------------- Problem 2 -------------------------
\miquestion TODO

%------------------------- Problem 3 -------------------------
\miquestion Review the definition of Thomae's function $t(x)$ from Section 4.1.
\begin{enumerate}[label=(\alph*)]
  \item Construct three different sequences $(x_{n})$, $(y_{n})$ and $(z_{n})$, each of which converges to 1 without using the number 1 as a term in the sequence.
  \item Now, compute $\lim t(x_{n})$, $\lim t(y_{n})$, and $\lim t(z_{n})$.
  \item Make an educated conjecture for $\lim_{x \to 1}t(x)$, and use Definition 4.2.1B to verify the claim. (Given $\epsilon > 0$, consider the set of points
  $\{x \in \textbf{R} : t(x) \geq \epsilon\}$). Argue that all the points in this set are isolated).
\end{enumerate}

\begin{solution}
  \begin{enumerate}[label=(\alph*)]
    \item A popular sequence that converges to 1 is $(x_{n}) = (\frac{n}{n+1})$.
    
    Then we can just add a sequence that converges to 0. By the Algebraic Limit Theorem
    the result will still converge to 1. So we can define $(y_{n}) = (x_{n}) + (\frac{1}{n})$ and $(z_{n}) = (x_{n}) + (\frac{1}{n^{2}})$.

    Some simple algebra will verify that none of the sequence has 1 as a term.

    \item Note that $t(\frac{p}{q}) = \frac{1}{q}$ for $\frac{p}{q}$ in its irreducible form. We need to make sure that the three sequences from
    part(a) are in their irreducible form.
    
    For $(x_{n}) = \frac{n}{n+1}$, the nume and denom are coprime (because they differ by 1) so they're already irreducible.
    Then $t(x_{n}) = (\frac{1}{n+1})$ and $\lim t(x_{n}) = 0$.

    For $(y_{n})$, we have $\frac{n}{n+1} + \frac{1}{n} = \frac{n^{2}+n+1}{n^{2}+n}$. Again the nume and denom differ by 1 so they're coprime. Then
    $t(y_{n}) = (\frac{1}{n^{2}+n})$ and $\lim t(y_{n}) = 0$.

    For $(z_{n})$, we have $\frac{n}{n+1} + \frac{1}{n^{2}} = \frac{n^3+n+1}{n^3+n^2}$. The nume and denom are coprime. To see this, rewrite the
    fraction as $\frac{n^3+n+1}{n^2(n+1)}$. Suppose they're not coprime, then there's a prime $p$ that divides both. Look at the denom. Since
    $n^2$ and $(n+1)$ are coprime, either $p \mid n$ or $p \mid (n+1)$. Both lead to contradiction.

    If $p \mid n$, then $(n^3+n+1) \mod p = 1$, contradicting the fact that $p$ divides the nume. If $p \mid (n+1)$, then
    $(n^3 + (n+1)) \mod p = n^3 + 0 \mod p$. Since $p \mid (n+1) \implies n \mod p = -1$, we have $n^3 \mod p = (-1)^3 = -1$.
    So $(n^3 + (n+1)) \mod p = -1 + 0 \mod p = -1$, again contradicting the fact that $p$ divides the nume.

    Then $t(z_{n}) = (\frac{1}{n^{3}+n^{2}})$ and $\lim t(z_{n}) = 0$.


    -------------------------------------\\
    Alternatively, we can choose an easier $(z_{n})$ like $\frac{2n+1}{2n}$ or $\frac{n^{2}+1}{n^{2}}$ that makes it instantly obvious that they're irreducible,
    avoiding the number theory mess.

    \item $\lim_{x \to 1}t(x) = 0$.  I initially didn't know how to prove it using the topological hint and instead used epsilon-delta.
    
    To use epsilon-delta, we need a way to determine a $V_{\delta}(1)$ neighborhood such that every (irreducible) rational between $(1-\delta, 1+\delta) \setminus 1$
    has a denominator greater than $n$ for some $n \in \textbf{N}$ i.e. $x \in V_{\delta}(1) \implies t(x) < \frac{1}{n}$. 
    Once we have this, it's easy to restrict $t(x)$ as small as we want.

    One intuitive guess is $(1-\frac{1}{n}, 1+\frac{1}{n})$. Is it true that every rational in-between that interval has denominator greater than $n$? Try some numbers,
    say $\delta = \frac{1}{4}$. Let's verify rationals with denom $< n$ are all outside $V_{\delta}(1)$
    \begin{enumerate}
      \item $\frac{1}{2}=(1-\frac{1}{2}) < (1-\frac{1}{4})$ (outside). $\frac{2}{2} = 1$ (outside, remember that we exclude 1). $\frac{3}{2} = (1+\frac{1}{2}) > (1+\frac{1}{4})$ (outside).
      \item $\frac{1}{3} < (1-\frac{1}{4})$ (outside). $\frac{2}{3}=(1-\frac{1}{3}) < (1-\frac{1}{4})$ (outside). $\frac{3}{3} = 1$ (outside). $\frac{4}{3} = (1+\frac{1}{3}) > (1+\frac{1}{4})$(outside).
    \end{enumerate}

    From the above example, it's easy to generalize the fact. Suppose $m < n$. Then $\frac{1}{m} > \frac{1}{n}$. Now, $\frac{m}{m}=1$ is outside. If we move left, we go to $(1-\frac{1}{m}) < (1-\frac{1}{n})$
    and that's already outside. It'll clearly still be outside if we move left even further. So any $\frac{m-k}{m}$ for $k \geq 0$ is outside the interval. A similar argument works for the right side.
    Combined, they imply that any $\frac{x}{m}$ with $m < n$ must be outside $V_{\delta}(1)$.

    The rest is simple. Given an $\epsilon > 0$, choose $n$ such that $\frac{1}{n} < \epsilon$ and set $\delta = \frac{1}{n}$. 
    Then we have $|x-1| < \delta \implies t(x) < \frac{1}{n} < \epsilon$.

    Note that $t(x) = 0$ for irrational $x$, so $t(x) < \epsilon$ in that case.
    
    TODO: the topological proof
  \end{enumerate}
\end{solution}

%------------------------- Problem 4 -------------------------
\miquestion TODO

%------------------------- Problem 5 -------------------------
\miquestion TODO

%------------------------- Problem 6 -------------------------
\miquestion True or false with short justifications.
\begin{enumerate}[label=(\alph*)]
  \item If a particular $\delta$ has been constructed as a suitable response to a particular $\epsilon$ challenge, then any smaller positive $\delta$ will also suffice.
  \item If $\lim_{x \to a}f(x) = L$ and $a$ happens to be in the domain of $f$, then $L = f(a)$.
  \item If $\lim_{x \to a}f(x) = L$, then $\lim_{x \to a}3[f(x)-2]^{2} = 3(L-2)^{2}$
  \item If $\lim_{x \to a}f(x) = 0$, then $\lim_{x \to a}f(x)g(x) = 0$ for any function $g$ (with domain equal to the domain of $f$).
\end{enumerate}

\begin{solution}
  \begin{enumerate}[label=(\alph*)]
    \item True. If $\delta_{1}$ is a suitable response it means $0 < |x-a| < \delta_{1} \implies |f(x)-L| < \epsilon$. Given $0 < \delta_{2} < \delta_{1}$, then
    $0 < |x-a| < \delta_{2} < \delta_{1}$.
    \item False. Let $\textbf{R}$ be the domain of $f$. Let $f(0) = 5$ and $f(x) = 0$ for $x \neq 0$. Then $\lim_{x \to 0} f(x) = 0$ but $f(0) = 5$.
    \item True. All the terms converge, including $(f(x)-2)$ that converges to $(L-2)$, so we can apply the Algebraic Limit Theorem.
    \item False. Let $\textbf{R} \setminus 0$ be the domain. Let $f(x) = x$ and $a=0$, then $\lim_{x \to 0}f(x) = 0$ as required. Let $g(x) = \frac{1}{x}$. 
    Then $f(x)g(x) = 1$ everywhere in the domain, but $\lim_{x \to 0}f(x)g(x) = 1$.
  \end{enumerate}
\end{solution}

%------------------------- Problem 7 -------------------------
\miquestion Let $g: A \to \textbf{R}$ and assume that $f$ is a bounded function on $A$ in the sense that there exists $M > 0$ satisfying $|f(x)| \leq M$
for all $x \in A$. Show that if $\lim_{x \to c}g(x) = 0$ then $\lim_{x \to c}g(x)f(x) = 0$ as well.

\begin{solution}
  We need to prove that for any $\epsilon > 0$ there exists $\delta$ such that $0 < |x-c| < \delta \implies |f(x)g(x)| < \epsilon$.

  Since $|f(x)| \leq M$, we have $|f(x)g(x)| \leq M|g(x)|$. Since $\lim_{x \to c}g(x) = 0$, we can pick $\delta$ such that
  $|g(x)| < \frac{\epsilon}{M} \implies M|g(x)| < M\frac{\epsilon}{M} = \epsilon$.
\end{solution}

%------------------------- Problem 8 -------------------------
\miquestion Compute each limit or state it does not exist. Use the tools developed in this section to justify each conclusion.
\begin{enumerate}[label=(\alph*)]
  \item $\lim_{x \to 2}\frac{|x-2|}{x-2}$
  \item $\lim_{x \to 7/4}\frac{|x-2|}{x-2}$
  \item $\lim_{x \to 0}(-1)^{\lceil 1/x \rceil}$
  \item $\lim_{x \to 0}\sqrt[3]{x}(-1)^{\lceil 1/x \rceil}$
\end{enumerate}

\begin{solution}
  \begin{enumerate}[label=(\alph*)]
    \item Diverges, as the function approaches 1 from $x > 2$ but -1 from $x < 2$. More formally, use Corollary 4.2.5
    (Divergence Criterion for Functional Limits).
    
    The sequence $(x_{n}) = (2 * \frac{n}{n+1})$ converges to 2 from the left and $\lim f(x_{n}) = -1$ as $f(x) = -1$ for all $x < 2$.
    Similarly, the sequence $(y_{n}) = (2 * \frac{n+1}{n})$ converges to 2 from the right and
    $\lim f(y_{n}) = 1$ as $f(x) = 1$ for all $x > 2$. We have $\lim x_{n} = \lim y_{n} = 2$ but $\lim f(x_{n}) \neq \lim f(y_{n})$, so the Divergence
    Criterion is satisfied.

    \item $\lim_{x \to 7/4}\frac{|x-2|}{x-2} = -1$. Because $\frac{7}{4} < 2$, there exists a neighborhood of $\frac{7}{4}$ where $x < 2$, so
    $\frac{|x-2|}{x-2} = \frac{-(x-2)}{x-2} = -1$.

    \item Diverges. Intuitively, the function should fluctuate rapidly between 1 and -1. Formally, use the Divergence Criterion.
    For everything below, $n \in \textbf{N}$ and $n > 0$. 
    
    The sequence $(1/2n)$ converges to 0 and $(-1)^{\lceil 1/\frac{1}{2n} \rceil} = (-1)^{2n} = 1$ for all $n$, so
    $\lim f(1/2n) = 1$.

    The sequence $(1/(2n+1))$ converges to 0 and $(-1)^{\lceil 1/\frac{1}{(2n+1)} \rceil} = (-1)^{2n+1} = -1$ for all $n$, so
    $\lim f(1/(2n+1)) = -1$.

    We have $\lim (1/2n) = \lim (1/(2n+1))$ but $\lim f(1/2n) \neq \lim f(1/(2n+1))$, so the Divergence Criterion is satisfied.

    \item $\lim_{x \to 0}\sqrt[3]{x}(-1)^{\lceil 1/x \rceil} = 0$. Use the result of exercise 4.2.7. $\lim_{x \to 0}\sqrt[3]{x}=0$ and
    $(-1)^{\lceil 1/x \rceil}$ is bounded (it's either -1 or 1), so their limit product is zero.
  \end{enumerate}
\end{solution}

%------------------------- Problem 9 -------------------------
\miquestion \textbf{(Infinite Limits)}. The statement $\lim_{x \to 0} 1/x^{2} = \infty$ certainly makes intuitive sense. To construct a rigorous definition in the
challenge-response style of Definition 4.2.1 for an infinite limit statement of this form, we replace the (arbitrarily small) $\epsilon > 0$ challenge with an
(arbitrarily large) $M > 0$ challenge:\\
\hspace*{1em} \textit{Definition}: $\lim_{x \to c}f(x)=\infty$ means that for all $M > 0$ we can find a $\delta > 0$ such that whenever $0 < |x-c| < \delta$, it
follows that $f(x) > M$.

\begin{enumerate}[label=(\alph*)]
  \item Show $\lim_{x \to 0}1/x^{2}=\infty$ in the sense described in the previous definition.
  \item Now, construct a definition for the statement $\lim_{x \to \infty}f(x) = L$. Show\\ $\lim_{x \to \infty}1/x=0$.
  \item What would a rigorous definition for $\lim_{x \to \infty}f(x)=\infty$ look like? Give an example of such a limit.
\end{enumerate}

\begin{solution}
  \begin{enumerate}[label=(\alph*)]
    \item To find a plausible delta: $\frac{1}{x^2} > M \implies |x| < \frac{1}{\sqrt{M}}$. 
    
    Let $\delta = \frac{1}{\sqrt{M}}$, then whenever $0 < |x| < \delta$, it follows that $\frac{1}{x^2} > \frac{1}{(\frac{1}{\sqrt{M}})^2} = M$.

    \item The definition is very similar to the definition of sequential limit. For all $\epsilon > 0$, there exists $M > 0$ such that $x > M \implies |f(x)-L| < \epsilon$.
    
    Pick $M = 1/\epsilon$. Then $x > M > 0 \implies \frac{1}{x} < \frac{1}{\frac{1}{\epsilon}} = \epsilon$.

    \item For all $L > 0$, there exists $M > 0$ such that $x > M \implies f(x) > L$.
    
    Example: $\lim_{x \to \infty}x = \infty$. Pick $M=L$. Then $x > (M=L) \implies (f(x)=x) > L$.
  \end{enumerate}
\end{solution}

%------------------------- Problem 10 -------------------------
\miquestion
\begin{enumerate}[label=(\alph*)]
  \item Give a proper definition in the style of Definition 4.2.1 for the right-hand and left-hand limit statements: $\lim_{x \to a^{+}}f(x) = L$ and
  $\lim_{x \to a^{-}}f(x) = M$.
  \item Prove that $\lim_{x \to a}f(x)=L$ if and only if both the right and left-hand limits equal $L$.
\end{enumerate}

\begin{solution}
  \begin{enumerate}[label=(\alph*)]
    \item Right limit: For all $\epsilon > 0$, there exists $\delta > 0$ such that $0 < x-a < \delta \implies |f(x)-L| < \epsilon$.\\
    Left limit: For all $\epsilon > 0$, there exists $\delta > 0$ such that $0 < a-x < \delta \implies |f(x)-L| < \epsilon$.\\
  
    \item We prove the sufficiency first: $\lim_{x \to a} f(x) = L \implies \\ \lim_{x \to a^{+}}f(x) = \lim_{x \to a^{-}}f(x) = L$.
    
    Pick $\delta$ such that $0 < |x-a| < \delta \implies |f(x)-L| < \epsilon$. For the left limit, we have $0 < a-x < \delta$ implies
    $0 < |x-a| < \delta$ which in turn implies $|f(x) - L| < \epsilon$. Similary for the right limit.

    For the other direction, we need to prove that $\lim_{x \to a^{+}}f(x) = \lim_{x \to a^{-}}f(x) = L \implies \lim_{x \to a} f(x) = L$.

    There exists $\delta_{L}$ such that $0 < a-x < \delta_{L} \implies |f(x) - L| < \epsilon$. Similarly, there exists $\delta_{R}$ such that
    $0 < x-a < \delta_{R} \implies |f(x) - L| < \epsilon$.

    Pick $\delta = \min(\delta_{L}, \delta_{R})$. Then $0 < |x-a| < \delta$ implies either $0 < a-x < \delta$ or $0 < x-a < \delta$, both of which
    imply $|f(x) - L| < \epsilon$.
  \end{enumerate}
\end{solution}

%------------------------- Problem 11 -------------------------
\miquestion \textbf{(Squeeze Theorem).} Let $f$, $g$, and $h$ satisfy $f(x) \leq g(x) \leq h(x)$ for all $x$ in some common domain $A$. If
$\lim_{x \to c}f(x) = \lim_{x \to c}h(x) = L$ at some limit point $c$ of $A$, show $\lim_{x \to c}g(x)=L$ as well.

\begin{solution}
  There exists $\delta$ such that $0 < |x-c| < \delta$ then $|f(x) - L| < \epsilon$ and $|h(x) - L| < \epsilon$. We need to prove that 
  using this $\delta$, $|g(x) - L| < \epsilon$ as well.

  For intuition, visualize the following. If $g(x) > L$, then $h(x)$ acts as an upper bound, preventing $g(x)$ from going over the upper boundary of $V_{\epsilon}(L)$.
  But if $g(x) < L$ then the upper bound is useless, as we can have $h(x)$ very close to $L-\epsilon$, and the distance $|h(x)-g(x)|$ pushes $g(x)$ below $L-\epsilon$.
  But in this case, $f(x)$ can act as a lower bound.

  Formally, we split into two cases:
  \begin{enumerate}
    \item $g(x) \geq L$. Since $h(x) \geq g(x)$ then $g(x)-L \leq h(x)-L$. But $|h(x)-L| < \epsilon$, so $g(x)-L \leq h(x)-L < \epsilon$.
    \item $g(x) \leq L$. Since $f(x) \leq g(x)$ then $L-g(x) \leq L-f(x)$. But $|f(x)-L| < \epsilon$, so $L-g(x) \leq L-f(x) < \epsilon$.
  \end{enumerate}

  Combining $g(x)-L < \epsilon$ and $L-g(x) < \epsilon$ gives $|g(x)-L| < \epsilon$.
\end{solution}

\end{questions}
\end{document}
